% SPDX-FileCopyrightText: 2026 Alexander S. Donahue
% SPDX-License-Identifier: CC-BY-4.0
\section{Stage 1: Exact Combinatorial Reduction (2D to 1D)}

\begin{definition}[Adjacent Pair Expectation $h_{\mathrm{adj}}(L)$]
\label{def:hadj}
\lean{MultinomialTV.hAdj, MultinomialTV.hAdj_eq_nested_sum}
\leanok
Defines $h_{\mathrm{adj}}(L) = \mathbb{E}[|X_2 - X_1|\,|X_3 - X_2|]$ under $\mathrm{Mult}(L; \frac13, \frac13, \frac13)$ on the discrete simplex and proves equivalence with the verbatim nested multinomial sum.
\end{definition}

\begin{theorem}[Exact 1D Reduction for $L = 3q$]
\label{thm:hadj-mod0}
\lean{MultinomialTV.hAdj_mod0}
\uses{def:hadj}
\leanok
Collapses the two-dimensional expectation $h_{\mathrm{adj}}(3q)$ identically into an exact rational combination of the central trinomial mass $p_{3q}(q)$ and the 1D sums $\mathrm{sumA}(q)$ and $\mathrm{sumB}(q)$.
\end{theorem}

\begin{theorem}[Exact 1D Reduction for $L = 3q + 1$]
\label{thm:hadj-mod1}
\lean{MultinomialTV.hAdj_mod1}
\uses{def:hadj}
\leanok
Collapses $h_{\mathrm{adj}}(3q+1)$ identically into an exact rational combination of $p_{3q}(q)$ and $\mathrm{sumA}(q)$.
\end{theorem}

\begin{theorem}[Exact 1D Reduction for $L = 3q - 1$]
\label{thm:hadj-mod2}
\lean{MultinomialTV.hAdj_mod2}
\uses{def:hadj}
\leanok
Collapses $h_{\mathrm{adj}}(3q-1)$ identically into an exact rational combination of $p_{3q}(q)$, $\mathrm{sumA}(q)$, $\mathrm{sumB}(q)$, and $\mathrm{sumC}(q)$.
\end{theorem}

\section{Stage 2: Quantitative Stirling \& Gaussian Euler--Maclaurin}

\begin{theorem}[Refined Central Trinomial Stirling Expansion]
\label{thm:stirling}
\lean{MultinomialTV.p3q_center_expansion}
\leanok
Establishes the quantitative $O(q^{-3})$ asymptotic expansion $p_{3q}(q) = \frac{\sqrt{3}}{2\pi q}\bigl(1 - \frac{2}{9q} + \frac{2}{81q^2} + O(q^{-3})\bigr)$ for all $q \ge 2$.
\end{theorem}

\begin{theorem}[Odd-Lattice Gaussian Euler--Maclaurin Bound]
\label{thm:gauss-em}
\lean{MultinomialTV.oddGauss_sum_approx}
\leanok
Proves quantitative Euler--Maclaurin summation accuracy for the odd-shifted Gaussian lattice sum driving the 1D local asymptotics.
\end{theorem}

\section{Stage 3: Local Asymptotics of 1D Residual Sums}

\begin{theorem}[Asymptotic Limit of $\mathrm{sumA}(q)$]
\label{thm:sumA}
\lean{MultinomialTV.tendsto_sumA}
\uses{thm:stirling, thm:gauss-em}
\leanok
Proves $q\bigl(\mathrm{sumA}(q) + \frac{1}{9} + \frac{4}{243q}\bigr) \to 0$ as $q \to \infty$.
\end{theorem}

\begin{theorem}[Asymptotic Limit of $\mathrm{sumB}(q)$]
\label{thm:sumB}
\lean{MultinomialTV.tendsto_sumB}
\uses{thm:stirling, thm:gauss-em}
\leanok
Proves $q\bigl(\mathrm{sumB}(q) - \frac{1}{3} + \frac{2}{27q}\bigr) \to 0$ as $q \to \infty$.
\end{theorem}

\begin{theorem}[Asymptotic Limit of $\mathrm{sumC}(q)$]
\label{thm:sumC}
\lean{MultinomialTV.tendsto_sumC}
\uses{thm:stirling, thm:gauss-em}
\leanok
Proves $q\bigl(\mathrm{sumC}(q) + \frac{2}{27q}\bigr) \to 0$ as $q \to \infty$.
\end{theorem}

\section{Stage 4: Residue Limits, Conjecture 7.7 \& Variance Closed Forms}

\begin{theorem}[Residue Limit $b_0$ for $L \equiv 0 \pmod{3}$]
\label{thm:mod0}
\lean{MultinomialTV.tendsto_hAdj_mod0, MultinomialTV.bRes_zero}
\uses{thm:hadj-mod0, thm:stirling, thm:sumA, thm:sumB}
\leanok
Proves $\lim_{q\to\infty} 3q\,(h_{\mathrm{adj}}(3q) - c_L(3q) - a_*) = \frac{32\sqrt{3}}{81\pi}$.
\end{theorem}

\begin{theorem}[Residue Limit $b_1$ for $L \equiv 1 \pmod{3}$]
\label{thm:mod1}
\lean{MultinomialTV.tendsto_hAdj_mod1, MultinomialTV.bRes_one}
\uses{thm:hadj-mod1, thm:stirling, thm:sumA}
\leanok
Proves $\lim_{q\to\infty} (3q+1)(h_{\mathrm{adj}}(3q+1) - c_L(3q+1) - a_*) = -\frac{22\sqrt{3}}{81\pi}$.
\end{theorem}

\begin{theorem}[Residue Limit $b_2$ for $L \equiv 2 \pmod{3}$]
\label{thm:mod2}
\lean{MultinomialTV.tendsto_hAdj_mod2, MultinomialTV.bRes_two}
\uses{thm:hadj-mod2, thm:stirling, thm:sumA, thm:sumB, thm:sumC}
\leanok
Proves $\lim_{q\to\infty} (3q+2)(h_{\mathrm{adj}}(3q+2) - c_L(3q+2) - a_*) = -\frac{22\sqrt{3}}{81\pi}$.
\end{theorem}

\begin{theorem}[Resolution of Conjecture 7.7]
\label{thm:conj77}
\lean{MultinomialTV.conjecture_7_7, MultinomialTV.conjecture_7_7_d0, MultinomialTV.conjecture_7_7_aTriangle, MultinomialTV.conjecture_7_7_gamma2, MultinomialTV.bRes_cosine_formula}
\uses{thm:mod0, thm:mod1, thm:mod2}
\leanok
Proves $d_0 = -\frac{4\sqrt{3}}{81\pi}$, $a_\triangle = \frac{1}{32} - \frac{4\sqrt{3}}{243}$, $\gamma_2 = -\frac{1}{32\pi} - \frac{8\sqrt{3}}{243\pi}$, and $b_r = d_0 + \frac{4\sqrt{3}}{9\pi}\cos(\frac{2\pi r}{3})$ for $r \in \{0,1,2\}$.
\end{theorem}

\begin{corollary}[Linear and Cyclic Variance Closed Forms $\gamma(n)$ and $\gamma_c(n)$]
\label{cor:gamma}
\lean{MultinomialTV.gamma_closed_form, MultinomialTV.gamma_cyclic_closed_form}
\uses{thm:conj77}
\leanok
Algebraic closed forms for $\gamma(n)$ (Part 2 Thm 7.6) and $\gamma_c(n)$ (Part 5 Thm 5.2) for all $n \ne 0$.
\end{corollary}

\section{Stage 5: Non-Vacuity Witnesses \& 6-Kernel Comparator Bridge}

\begin{proposition}[Non-Vacuity \& Adversarial Refutation Witnesses]
\label{prop:nonvac}
\lean{MultinomialTV.NonVacuity.hAdj_zero, MultinomialTV.NonVacuity.hAdj_one, MultinomialTV.NonVacuity.hAdj_two, MultinomialTV.NonVacuity.hAdj_three, MultinomialTV.NonVacuity.hAdj_four, MultinomialTV.NonVacuity.hAdj_five, MultinomialTV.NonVacuity.bRes_zero_pos, MultinomialTV.NonVacuity.bRes_one_neg, MultinomialTV.NonVacuity.bRes_zero_ne_one, MultinomialTV.NonVacuity.d0_neg, MultinomialTV.NonVacuity.aTriangle_lt_one_div_32, MultinomialTV.NonVacuity.aTriangle_pos, MultinomialTV.NonVacuity.gamma2_neg, MultinomialTV.NonVacuity.adversarial_aTriangle_ne_one_div_32, MultinomialTV.NonVacuity.adversarial_d0_ne_zero, MultinomialTV.NonVacuity.adversarial_residue_limits_not_constant, MultinomialTV.NonVacuity.gamma_closed_form_witness_n5, MultinomialTV.NonVacuity.gamma_cyclic_closed_form_witness_n5}
\uses{def:hadj, thm:conj77, cor:gamma}
\leanok
18 concrete small-$L$, strict-inequality, adversarial false-constant refutation, and $n=5$ inhabitation theorems ruling out vacuous definitions.
\end{proposition}

\begin{theorem}[Isolated 6-Kernel Comparator Targets]
\label{thm:comparator}
\lean{MultinomialTV.ComparatorTarget.target_tendsto_hAdj_mod0, MultinomialTV.ComparatorTarget.target_tendsto_hAdj_mod1, MultinomialTV.ComparatorTarget.target_tendsto_hAdj_mod2, MultinomialTV.ComparatorTarget.target_conjecture_7_7, MultinomialTV.ComparatorTarget.target_gamma_closed_form, MultinomialTV.ComparatorTarget.target_gamma_cyclic_closed_form}
\uses{thm:mod0, thm:mod1, thm:mod2, thm:conj77, cor:gamma}
\leanok
Mathlib-only isolated specification (\texttt{Challenge.lean}, 122 lines, Lean 4.35.0 \texttt{module} header) matched against sorry-free \texttt{Solution.lean} by 6 independent kernels via \texttt{lake comparator --config comparator.json --paranoid}.
\end{theorem}
